The ElGamal public-key encryption scheme \cite{C:ElGamal84} in groups of prime order has been proven IND-CPA secure under the DDH assumption \cite{ElGamalSec,C:CraSho98,NaoRei04}.
Let $\groupgen$ be an algorithm that takes as input the security parameter $\secparam$ and outputs a group description $(\GG, p, g)$ where $\GG$ is a cyclic group of order $p$ with generator $g$.
The message space $\msgspace$ is the generated group itself.
A description of the scheme can be seen in \hme{add when added}.
% TODO: NORMALIZE MACROS
% \begin{figure}[h]
% 	\begin{pcvstack}[boxed, center, space=1em]
% 		\procedure[linenumbering]{$\setup(\secparam)$}{
% 			\pcreturn \groupgen(\secparam)
% 		}
%
% 		\begin{pchstack}
%
% 			\procedure[linenumbering]{$\kgen(\pkepp)$}{
% 				\pcparse (\GG, g, p) = \pkepp \\
% 				x \sample \ZZ_{p} \\
% 				X \leftarrow g^x \\
% 				\pk \leftarrow(p, g, X) \\
% 				\sk \leftarrow(p, g, x) \\
% 				\pcreturn (\sk, \pk)
% 			}
% 			\procedure[linenumbering]{$\enc(\pkepp, \pk, M)$}{
% 				\pcparse (\group, \ggen, \egprime) = \pkepp \\
% 				y \sample \Z_{\egprime} \\
% 				Y \leftarrow \ggen^y \\
% 				T \leftarrow X^y \\
% 				W \leftarrow T M \\
% 				\pcreturn (Y, W)
% 			}
% 			\procedure[linenumbering]{$\dec(\pkepp, sk, ct)$}{
% 				(Y, W) \leftarrow ct  \\
% 				T \leftarrow Y^x \\
% 				M \leftarrow W T^{-1} \\
% 				\pcreturn M
% 			}
% 		\end{pchstack}
% 	\end{pcvstack}
% 	\caption{The ElGamal Encryption Scheme}
% 	\label{fig:regelg}
% \end{figure}
